\documentclass[12pt,a4paper]{article}
\usepackage{ugmath}
\usepackage{float}
\usepackage{placeins}
\usepackage{array}
\usepackage[labelfont=bf,labelsep=space]{caption}
\newcommand{\paren}[1]{\left( #1 \right)}
\newcommand{\alumno}{Ricardo León Martínez}
\newcommand{\materia}{Elementos de Probabilidad y Estadistica}
\newcommand{\profesor}{Claudia Reynoso Alcántara}
\newcommand{\tarea}{Tarea 5}
\newcommand{\fecha}{25/3/2026}

\begin{document}

    \begin{center}
        {\large\textbf{UNIVERSIDAD DE GUANAJUATO}}\\[0.3cm]
        {\normalsize\textbf{DIVISIÓN DE CIENCIAS NATURALES Y EXACTAS}}\\
        {\normalsize\textbf{CAMPUS GUANAJUATO}}\\[1cm]

        {\Large\textbf{\tarea\ (\materia)}}\\[1cm]
    \end{center}

    \noindent
    \textbf{Nombre:} \alumno 
    \hfill 
    \textbf{Fecha:} \fecha 
    \hfill 
    \textbf{Calificación:} \rule{3cm}{0.4pt}

    \vspace{0.3cm}

    \begin{mdframed}[style=mdpinkbox,frametitle={Ejercicio 1}]
    Sean $X_{1},X_{2},\ldots,X_{n}$ variables aleatorias sobre un espacio medible $(\Omega,\mathcal{F})$ para $n\geq 2$.
        \begin{enumerate}
            \item[(a)] Demuestre a partir de la definición de variable aleatoria que
            \[
            \max\{X_{1},X_{2}\}
            \quad \text{y} \quad
            \min\{X_{1},X_{2}\}
            \]
            son variables aleatorias.

            \item[(b)] Demuestre que
            \[
            Y=\max\{X_{1},\ldots,X_{n}\}
            \quad \text{y} \quad
            Z=\min\{X_{1},\ldots,X_{n}\}
            \]
            son variables aleatorias.

            \textbf{Sugerencia.} Encuentre relaciones entre
            \[
            \max\{a,b,c\}
            \quad \text{y} \quad
            \max\{\max\{a,b\},c\},
            \]
            y también entre
            \[
            \min\{a,b,c\}
            \quad \text{y} \quad
            \min\{\min\{a,b\},c\}.
            \]
        \end{enumerate}
    \end{mdframed}

    \begin{proof}
        \begin{enumerate}
            \item[(a)] Definamos
            \begin{equation*}
                Y(\omega)=\max\{X_{1}(\omega),X_{2}(\omega)\}.
            \end{equation*}
            Observemos que
            \begin{equation*}
                Y(\omega)\leq t
            \end{equation*}
            si y solo si
            \begin{equation*}
                X_{1}(\omega)\leq t\quad\text{ y }\quad X_{2}(\omega)\leq t.
            \end{equation*}
            Por tanto,
            \begin{equation*}
                \{\omega: Y(\omega)\leq t\}=\{\omega:X_{1}(\omega)\leq t\}\cap\{\omega:X_{2}(\omega)\leq t\}
            \end{equation*}
            Como $X_{1}$ y $X_{2}$ son variables aleatorias,
            \begin{equation*}
                \{\omega:X_{1}(\omega)\leq t\}\in\FF,\qquad \{\omega:X_{2}(\omega)\leq t\}\in\FF.
            \end{equation*}
            Dado que $\FF$ es una $\sigma$-algebra y es cerrada bajo intersecciones finitas, se concluye que
            \begin{equation*}
                \{\omega: Y(\omega)\leq t\}\in\FF
            \end{equation*}
            Luego, $\max\{X_{1},X_{2}\}$ es una variable aleatoria. Ahora definamos
            \begin{equation*}
                Z(\omega)=\min\{X_{1}(\omega),X_{2}(\omega)\}.
            \end{equation*}
            Observemos que
            \begin{equation*}
                Z(\omega)\leq t
            \end{equation*}
            si y solo si
            \begin{equation*}
                X_{1}(\omega)\leq t\quad\text{ y }\quad X_{2}(\omega)\leq t.
            \end{equation*}
            Por tanto,
            \begin{equation*}
                \{\omega: Z(\omega)\leq t\}=\{\omega:X_{1}(\omega)\leq t\}\cap\{\omega:X_{2}(\omega)\leq t\}
            \end{equation*}
            Como $X_{1}$ y $X_{2}$ son variables aleatorias,
            \begin{equation*}
                \{\omega:X_{1}(\omega)\leq t\}\in\FF,\qquad \{\omega:X_{2}(\omega)\leq t\}\in\FF.
            \end{equation*}
            Dado que $\FF$ es una $\sigma$-algebra y es cerrada bajo intersecciones finitas, se concluye que
            \begin{equation*}
                \{\omega: Z(\omega)\leq t\}\in\FF
            \end{equation*}
            Luego, $\min\{X_{1},X_{2}\}$ es una variable aleatoria.
            \item[(b)] Procederemos por inducción sobre $n$. Para $n=2$, el resultado ya fue demostrado en el inciso anterior.
            Supongamos que
            \begin{equation*}
                \max\{X_{1},\dots,X_{n}\}\quad\text{ y }\quad\min\{X_{1},\dots,X_{n}\}
            \end{equation*}
            son variables aleatorias. Definamos
            \begin{equation*}
                M_{n}=\max\{X_{1},\dots,X_{n}\}\quad\text{ y }\quad m_{n}=\min\{X_{1},\dots,X_{n}\}.
            \end{equation*}
            Entonces ,
            \begin{equation*}
                \max\{X_{1},\dots,X_{n},X_{n+1}\}=\max\{M_{n},X_{n+1}\}\quad\text{ y }\quad\min\{X_{1},\dots,X_{n},X_{n+1}\}=\min\{m_{n},X_{n+1}\}
            \end{equation*}
            Por hipotesis inductiva, $M_{n}$ y $m_{n}$ son variables aleatorias. Ademas $X_{n+1}$ tambien es una variable
            aleatoria. Aplicando (a), concluimos que
            \begin{equation*}
                \max\{M_{n},X_{n+1}\}\quad\text{ y }\quad\min\{m_{n},X_{n+1}\}
            \end{equation*}
            son variables aleatorias. Así $\max\{X_{1},\dots,X_{n+1}\}$ y $\min\{X_{1},\dots,X_{n+1}\}$ son variables
            aleatorias. Por inducción el resultado es valido par atodo $n\geq2$.
        \end{enumerate}
    \end{proof}

    \begin{mdframed}[style=mdpinkbox,frametitle={Ejercicio 2}]
        Sean $X,Y$ variables aleatorias sobre $(\Omega,\mathcal{F})$. Demuestre que la función
        \[
        Z:\Omega\to\mathbb{R}^{2}
        \]
        dada por
        \[
        Z(\omega)=(X(\omega),Y(\omega))
        \]
        para todo $\omega\in\Omega$ es una variable aleatoria.

        ¿Se cumple lo análogo para la función
        \[
        \omega\longmapsto (X_{1}(\omega),X_{2}(\omega),\ldots,X_{n}(\omega))
        \]
        con $X_{1},X_{2},\ldots,X_{n}$ variables aleatorias?
    \end{mdframed}

    \begin{proof}
        Definamos
        \begin{equation*}
            Z(\o)=(X(\o),Y(\o))
        \end{equation*}
        Para ver si $Z$ es una variable aleatoria en $\R^{2}$, basta verificar si para todo $a,b\in\R$,
        \begin{equation*}
            \{\o\in \Omega:X(\o)\leq a, Y(\o)\leq b\}\in\FF.
        \end{equation*}
        Pero
        \begin{equation*}
            \{\o\in \Omega:X(\o)\leq a, Y(\o)\leq b\}=\{\o\in \Omega:X(\o)\leq a\}\cap\{\o\in\Omega:Y(\o)\leq b\}
        \end{equation*}
        Como $X$ y $Y$ son variables aleatorias,
        \begin{equation*}
            \{\omega:X(\omega)\leq a\}\in\FF,\qquad\{\omega:Y(\omega)\leq b\}\in\FF.
        \end{equation*}
        Ademas, $\FF$ es una $\sigma$-algebra, luego es cerrada bajo intersecciones finitas. Por tanto
        \begin{equation*}
            \{\o\in \Omega:X(\o)\leq a, Y(\o)\leq b\}\in\FF.
        \end{equation*}
        Concluimos que
        \begin{equation*}
            Z(\omega)=(X(\omega),Y(\omega))
        \end{equation*}
        es una variable aleatoria.\\
        Si se cumple el resultado para $n$ variables aleatorias. Si
        \begin{equation*}
            X_{1},X_{2},\dots,X_{n}
        \end{equation*}
        son variables aleatorias, definimos
        \begin{equation*}
            Z(\omega)=(X_{1}(\omega),X_{2}(\omega),\dots,X_{n}(\omega)).
        \end{equation*}
        Entonces, para culquiera $a_{1},a_{2},\dots,a_{n}\in\R$,
        \begin{align*}
            \{\omega\in\Omega:X_{1}(\omega)\leq a_{1},\dots,X_{n}(\omega)\leq a_{n}\}\\
            =\bigcap_{i=1}^{n}\{\omega:X_{i}(\omega)\leq a_{i}\}.
        \end{align*}
        Cada conjunto $\{\omega:X_{i}\leq a_{i}\}$ pertenece a $\FF$ porque $X_{i}$ es una variable aleatoria.
        Como $\FF$ es cerrada bajo intersecciones finitas,
        \begin{equation*}
            \bigcap_{i=1}^{n}\{\omega:X_{i}(\omega)\leq a_{i}\}\in\FF.
        \end{equation*}
        Por lo tanto,
        \begin{equation*}
            \omega\longmapsto(X_{1}(\omega),X_{2}(\omega),\dots,X_{n}(\omega))
        \end{equation*}
        es una variable aleatoria.
    \end{proof}

    \begin{mdframed}[style=mdpinkbox,frametitle={Ejercicio 3}]
        Sea $(\Omega,\mathcal{F})$ un espacio medible y fijemos un evento $A\in\mathcal{F}$. Sean $X,Y$ variables aleatorias sobre $(\Omega,\mathcal{F})$.

        Demuestre que la función dada por
        \[
        Z:\Omega\to\mathbb{R}
        \]
        \[
        Z(\omega)=
        \begin{cases}
            X(\omega), & \text{si } \omega\in A,\\
            Y(\omega), & \text{si } \omega\notin A,
        \end{cases}
        \]
        para todo $\omega\in\Omega$, es una variable aleatoria.
    \end{mdframed}

    \begin{proof}
        Sea $t\in\R$. Por definición de $Z$,
        \begin{equation*}
            Z(\o)\leq t
        \end{equation*}
        si y solo si ocurre una de las siguientes posibilidades
        \begin{equation*}
            \omega\in A\text{ y }X(\omega)\leq t,\qquad\omega\notin A\text{ y }Y(\omega)\leq t.
        \end{equation*}
        Por tanto,
        \begin{equation*}
            \{\o:Z(\o)\leq t\}=\paren{A\cap\{\omega:X(\omega)\leq t\}}\cup\paren{A^{c}\cap\{\o:Y(\o)\leq t\}}.
        \end{equation*}
        Como $X$ es variable aleatoria,
        \begin{equation*}
            \{\o:X(\o)\leq t\}\in\FF.
        \end{equation*}
        De igual forma, como $Y$ es variable aleatoria,
        \begin{equation*}
            \{\o:Y(\o)\leq t\}\in\FF.
        \end{equation*}
        Además, $A\in\FF$ y como $\FF$ es una $\sigma$-álgebra, $A^{c}\in\FF$. Dado que $\FF$ es cerrada bajo
        intersecciones finitas y uniones finitas
        \begin{equation*}
            A\cap\{\omega:X(\omega)\leq t\}\in\FF,\quad\text{ y }\quad A^{c}\cap\{\o:Y(\o)\leq t\}.
        \end{equation*}
        Así,
        \begin{equation*}
            \{\o:Z(\o)\leq t\}\in\FF.
        \end{equation*}
        Como esto ocurre para todo $t\in\R$, concluimos que $Z$ es una variable aleatoria.
    \end{proof}

    \begin{mdframed}[style=mdpinkbox,frametitle={Ejercicio 4}]
        Sea $(\Omega,\mathcal{F})$ el espacio medible donde
        \[
        \Omega=\{1,2,3,4,5,6\}
        \]
        y
        \[
        \mathcal{F}=\{\varnothing,\{2,4,6\},\{1,3,5\},\Omega\}.
        \]

        Consideremos $U,V,W:\Omega\to\mathbb{R}$ funciones definidas por
        \[
        U(\omega)=5\omega+32,
        \]
        \[
        V(\omega)=
        \begin{cases}
            1, & \text{si } \omega \text{ es par},\\
            0, & \text{si } \omega \text{ es impar},
        \end{cases}
        \]
        y
        \[
        W(\omega)=\omega^{2}.
        \]

        ¿Cuáles de estas funciones son variables aleatorias con respecto a $(\Omega,\mathcal{F})$?

        ¿Existe alguna $\sigma$-álgebra $\mathcal{G}$ tal que $U,V,W$ son variables aleatorias respecto a $(\Omega,\mathcal{G})$?
    \end{mdframed}

    \begin{solucion}
        Analisis de $U$.\\
        Calculemos los valores de $U$
        \begin{align*}
            U(1)=37,\\
            U(2)=42,\\
            U(3)=47,\\
            U(4)=52,\\
            U(5)=57,\\
            U(6)=62.
        \end{align*}
        Tomemos $t=40$. Entonces
        \begin{equation*}
            \{\o:U(\o)\leq40\}=\{1\}.
        \end{equation*}
        Pero $\{1\}\notin\FF$. Por lo tanto, $U$ no es variable aleatoria respecto a $(\Omega,\FF)$.\\
        Analisis de $V$.\\
        Tomemos cualquier $t\in\R$. Si $t\lt0$,
        \begin{equation*}
            \{\o:V(\o)\leq t\}=\varnothing.
        \end{equation*}
        Si $0\leq t\lt1$,
        \begin{equation*}
            \{\o:V(\o)\leq t\}=\{1,3,5\}.
        \end{equation*}
        Si $t\geq1$,
        \begin{equation*}
            \{\o:V(\o)\leq t\}=\Omega.
        \end{equation*}
        Todos estos conjuntos pertenecen a $\FF$. Luego $V$ es una variable aleatoria respecto a $(\Omega,\FF)$.\\
        Analisis de $W$.\\
        Calculemos
        \begin{align*}
            W(1)=1,\\
            W(2)=4,\\
            W(3)=9,\\
            W(4)=16,\\
            W(5)=25,\\
            w(6)=36.
        \end{align*}
        Tomemos $t=10$. Entonces
        \begin{equation*}
            \{\o:W(\o)\leq10\}=\{1,2,3\}.
        \end{equation*}
        Pero $\{1,2,3\}\notin\FF$. Así, $W$ no es variable aleatoria respecto a $(\Omega,\FF)$.\\
        Sí existe una $\sigma$-algebra $\mathcal{G}$ tal que $U,V,W$ sean variables aleatorias. Por ejemplo
        se puede tomar
        \begin{equation*}
            \mathcal{G}=\mathcal{P}(\Omega),
        \end{equation*}
        la $\sigma$-algebra potencia de $\Omega$, es decir, el conjunto de todos los subconjuntos de $\Omega$.
        Así, para cuialquier función $X:\Omega\to\R$ y cualqueir $t\in\R$,
        \begin{equation*}
            \{\o:X(\o)\leq t\}\subseteq\Omega.
        \end{equation*}
        Como todos los subconjuntos de $\Omega$ pertenecen a $\mathcal{P}(\Omega)$, se concluye que toda función
        $X:\Omega\to\R$ es variable aleatoria respecto a $(\Omega,\mathcal{P}(\Omega))$. Por tanto $U,V,W$ son
        variables aleatorias respecto a $(\Omega,\mathcal{P}(\Omega))$.
    \end{solucion}

    \begin{mdframed}[style=mdpinkbox,frametitle={Ejercicio 5}]
    Sea $X$ una variable aleatoria con función de probabilidad, es decir,
    \[
    P(X=r_{i})=p_{i},
    \]
    dada por la siguiente tabla:

    \[
    \begin{array}{c|cccccccc}
        r_{i} & -2 & -1 & 0 & 1 & 2 & 3 & 4 & 5\\
        \hline
        p_{i} & 0.1 & 0.2 & 0.15 & 0.2 & 0.1 & 0.15 & 0.05 & 0.05
    \end{array}
    \]

    Calcule la probabilidad de cada uno de los siguientes eventos.

        \begin{enumerate}
            \item[(a)] $X$ es negativa.
            \item[(b)] $X$ es par.
            \item[(c)] $X$ toma valores en el intervalo $[1,5]$.
            \item[(d)] $P(X=-2\mid X\leq 0)$.
            \item[(e)] $P(X\geq 2\mid X>0)$.
        \end{enumerate}
    \end{mdframed}

    \begin{solucion}
        \begin{enumerate}
            \item[(a)] $X$ es negativa.\\
            El evento es
            \begin{equation*}
                \{X\lt0\}=\{-2,-1\}.
            \end{equation*}
            Por tanto,
            \begin{equation*}
                P(X\lt0)=P(X=-2)+P(X=-1).
            \end{equation*}
            Entonces,
            \begin{equation*}
                P(X\lt0)=0.1+0.2=0.3.
            \end{equation*}
            \begin{equation*}
                P(X\lt0)=0.3
            \end{equation*}
            \item[(b)] $X$ es par\\
            Los valores pares son -2,0,2,4. Así
            \begin{align*}
                P(X\text{ es par})&=P(X=-2)+P(X=0)+P(X=2)+P(X=4)\\
                &=0.1+0.15+0.1+0.005\\
                &=0.4.
            \end{align*}
            \begin{equation*}
                P(X\text{ es par})=0.4.
            \end{equation*}
            \item[(c)] $X$ toma los valores en el intervalo $[1,5]$\\
            El evento es $\{1,2,3,4,5\}$. Entonces
            \begin{align*}
                P(1\leq X\leq 5)&=0.2+0.1+0.15+0.05+0.05\\
                &=0.55.
            \end{align*}
            \begin{equation*}
                P(1\leq X\leq 5)=0.55
            \end{equation*}
            \item[(d)] $P(X=-2\mid X\leq0)$\\
            Tomemos $A=\{X=-2\}$, $B=\{X\leq0\}$. Como $A\subseteq B$, se tiene $A\cap B=A$. Por tanto,
            \begin{equation*}
                P(X=-2\mid X\leq0)=\frac{P(X=-2)}{P(X\leq0)}.
            \end{equation*}
            Ahora,
            \begin{equation*}
                P(X\leq0)=0.1+0.2+0.15=0.45.
            \end{equation*}
            Entonces,
            \begin{equation*}
                P(X=-2\mid X\leq0)=\frac{0.1}{0.45}=\frac{2}{9}.
            \end{equation*}
            \begin{equation*}
                P(X=-2\mid X\leq0)=\frac{2}{9}
            \end{equation*}
            \item[(e)] Tomemos $A=\{X\geq2\}=\{2,3,4,5\}$, y $B=\{X\gt0\}=\{1,2,3,4,5\}$. Entonces $A\subseteq B$,
            y por lo tanto
            \begin{equation*}
                P(X\geq2\mid X\gt0)=\frac{P(X\geq2)}{X\gt0}.
            \end{equation*}
            Calculemos
            \begin{equation*}
                P(X\geq2)=0.1+0.15+0.05+0.05=0.35,
            \end{equation*}
            y
            \begin{equation*}
                P(X\gt0)=0.2+0.1+0.15+0.05+0.05=0.55.
            \end{equation*}
            Así,
            \begin{equation*}
                P(X\geq2\mid X\gt0)=\frac{0.35}{0.35}=\frac{7}{11}.
            \end{equation*}
            \begin{equation*}
                P(X\geq2\mid X\gt0)=\frac{7}{11}
            \end{equation*}
        \end{enumerate}
    \end{solucion}

    \begin{mdframed}[style=mdpinkbox,frametitle={Ejercicio 6}]
        Sea $X$ una variable distribuida uniformemente en el conjunto
        \[
        \{1,2,\ldots,50\},
        \]
        esto es,
        \[
        P(X=k)=\frac{1}{50}
        \]
        para $k=1,2,\ldots,50$.

        Encuentre las siguientes probabilidades.

        \begin{enumerate}
            \item[(a)] $P(X\geq 15)$.
            \item[(b)] $P(2.5<X\leq 43.2)$.
            \item[(c)] $P(X>20\mid X>10)$.
            \item[(d)] $P\left(X\leq \frac{435}{6}\,\middle|\,X>15\right)$.
        \end{enumerate}
    \end{mdframed}

    \begin{solucion}
        Como todos los valores equiprobables, la probabilidad de un evento es
        \begin{equation*}
            P(A)=\frac{\#A}{50},
        \end{equation*}
        donde $\#A$ es el número de elementos del conjunto $A$.
        \begin{enumerate}
            \item[(a)] Los valores posibles son $15,16,\dots,50$. La cantidad de enteros es
            \begin{equation*}
                50-15+1=36.
            \end{equation*}
            Por tanto,
            \begin{equation*}
                P(X\geq15)=\frac{36}{50}=\frac{18}{25}.
            \end{equation*}
            \begin{equation*}
                P(X\geq15)=\frac{18}{25}
            \end{equation*}
            \item[(b)] Como $X$ toma los valores enteros, la condición equivale a
            \begin{equation*}
                3\leq X\leq43.
            \end{equation*}
            Los valores posibles son $3,4,\dots,43$. La cantidad de enteros es
            \begin{equation*}
                43-3+1=41.
            \end{equation*}
            Entonces,
            \begin{equation*}
                P(2.5\lt X\leq43.2)=\frac{41}{50}.
            \end{equation*}
            \item[(c)] Tomemos $A=\{X\gt20\}$, $B=\{X\gt10\}$. Como $A\subseteq B$, se tiene $A\cap B=A$.
            Entonces,
            \begin{equation*}
                P(X\gt20\mid X\gt10)=\frac{P(X\gt20)}{X\gt10}.
            \end{equation*}
            Ahora $X\gt20$ corresponde a $21,22,\dots,50$, que son
            \begin{equation*}
                50-21+1=30
            \end{equation*}
            valores. Asimismo, $X\gt0$ corresponde a $11,12,\dots,50$, que son
            \begin{equation*}
                50-11+1=40
            \end{equation*}
            valores. Por tanto
            \begin{equation*}
                P(X\gt20\mid X\gt10)=\frac{\frac{30}{50}}{\frac{40}{50}}=\frac{30}{40}=\frac{3}{4}.
            \end{equation*}
            \begin{equation*}
                P(X\gt20\mid X\gt10)=\frac{3}{4}
            \end{equation*}
            \item[(d)]
            \begin{equation*}
                P\paren{X\leq\frac{435}{6}\mid X\gt15}
            \end{equation*}
            Primero observemos que
            \begin{equation*}
                \frac{435}{6}=72.5.
            \end{equation*}
            Como $X\in\{1,\dots,50\}$, siempre se cumple
            \begin{equation*}
                X\leq72.5.
            \end{equation*}
            Por tanto
            \begin{equation*}
                \left\{X\leq\frac{435}{6}\right\}=\Omega.
            \end{equation*}
            Así,
            \begin{equation*}
                P\paren{X\leq\frac{435}{6}\mid X\gt15}=P(\Omega\mid X\gt15).
            \end{equation*}
            Y sabemos que
            \begin{equation*}
                P(\Omega\mid B)=1
            \end{equation*}
            para cualquier evento $B$ con probabilidad positivia. Luego,
            \begin{equation*}
                P(X\leq\frac{435}{6}\mid X\gt15)=1.
            \end{equation*}
        \end{enumerate}
    \end{solucion}

    \begin{mdframed}[style=mdpinkbox,frametitle={Ejercicio 7}]
        Sea $X$ una variable aleatoria con función de probabilidad, es decir,
        \[
        P(X=r_{i})=p_{i},
        \]
        dada por la siguiente tabla:
        \[
        \begin{array}{c|cccccc}
            r_{i} & -2 & -1 & 0 & 1 & 2 & 3\\
            \hline
            p_{i} & 0.1 & 0.2 & 0.15 & 0.25 & 0.15 & 0.15
        \end{array}
        \]
        Sea $Y$ la variable aleatoria definida por
        \[
        Y=X^{2}.
        \]
        Halle la función de probabilidad de $Y$.
        Calcule el valor de la función de distribución de $X$ y de $Y$ en los puntos
        \[
        1,\quad \frac{3}{4},\quad \pi-3.
        \]
    \end{mdframed}

    \begin{proof}
        Los valores posibles de $Y$ son 0,1,4,9. Ahora agrupamos probabilidades. Para $Y=0$
        \begin{equation*}
            Y=0\ssi X=0.
        \end{equation*}
        Entonces,
        \begin{equation*}
            P(Y=0)=P(X=0)=0.15.
        \end{equation*}
        Para $Y=1$
        \begin{equation*}
            Y=1\ssi X=\pm1.
        \end{equation*}
        Así,
        \begin{align*}
            P(Y=1)&=P(X=-1)+P(X=1)\\
            &=0.2+0.25\\
            &=0.45.
        \end{align*}
        Para $Y=4$
        \begin{equation*}
            Y=4\ssi X=\pm2.
        \end{equation*}
        Entonces,
        \begin{align*}
            P(Y=4)&=P(X=-2)+P(X=2)\\
            &=0.1+0.15\\
            &=0.25.
        \end{align*}
        Para $Y=9$
        \begin{equation*}
            Y=9\ssi X=3.
        \end{equation*}
        Por lo tanto,
        \begin{equation*}
            P(Y=9)=0.15.
        \end{equation*}
        La función de probabilidad es
            \begin{equation*}
                \begin{array}{c|cccc}
                y & 0 & 1 & 4 & 9\\
                \hline
                P(Y=y) & 0.15 & 0.45 & 0.25 & 0.15
            \end{array}
        \end{equation*}
        Ahora recordemos que
        \begin{equation*}
            F_{X}(t)=P(X\leq t).
        \end{equation*}
        Valor en 1
        \begin{equation*}
            F_{X}(1)=P(X\leq1).
        \end{equation*}
        Entonces,
        \begin{align*}
            F_{X}(1)&=0.1+0.2+0.15+0.25\\
            &=0.7.
        \end{align*}
        \begin{equation*}
            F_{X}(1)=0.7
        \end{equation*}
        Valor en $\frac{3}{4}$. Como $\frac{3}{4}=0.75$. Por tanto,
        \begin{align*}
            F_{X}\paren{\frac{3}{4}}&=0.1+0.2+0.15\\
            &=0.45
        \end{align*}
        \begin{equation*}
            F_{X}\paren{\frac{3}{4}}=0.45
        \end{equation*}
        Valor en $\pi-3$. Observemos que
        \begin{equation*}
            \pi-3\approx0.14159.
        \end{equation*}
        Los valores de $X$ menores o iguales a $\pi-3$ son -2,-1,0. Entonces,
        \begin{align*}
            F_{X}(\pi-3)&=0.1+0.2+0.15\\
            &=0.45.
        \end{align*}
        \begin{equation*}
            F_{X}(\pi-3)=0.45
        \end{equation*}
        Ahora recordemos nuevamente que
        \begin{equation*}
            F_{Y}(t)=P(Y\leq t).
        \end{equation*}
        Los valores posibles de $Y$ son 0,1,4,9. Valor en 1
        \begin{equation*}
            F_{Y}(1)=P(Y\leq1).
        \end{equation*}
        Esto incluye $Y=0.1$. Entonces,
        \begin{align*}
            F_{Y}(1)&=0.15+0.45\\
            &=0.6.
        \end{align*}
        \begin{equation*}
            F_{Y}(1)=0.6
        \end{equation*}
        Valor en $\frac{3}{4}$. Como $\frac{3}{4}\lt1$, el único valor de $Y$ menor o igual que $\frac{3}{4}$
        es 0. Por tanto,
        \begin{equation*}
            F_{Y}\paren{\frac{3}{4}}=0.15.
        \end{equation*}
        Valor en $\pi-3$. Como
        \begin{equation*}
            \pi-3\approx0.14159\lt1,
        \end{equation*}
        nuevamente el único valor posible $Y$ menor o igual que $\pi-3$ es 0. Así,
        \begin{equation*}
            F_{Y}(\pi-3)=0.15.
        \end{equation*}
    \end{proof}

    \begin{mdframed}[style=mdpinkbox,frametitle={Ejercicio 8}]
        Sea $X$ una variable aleatoria con función de distribución $F_{X}$ dada por
        \[
        F_{X}(r)=
        \begin{cases}
            0, & \text{para } r<0,\\[0.3cm]
            \dfrac{1}{2}, & \text{para } 0\leq r<\dfrac{1}{4},\\[0.4cm]
            \dfrac{3}{4}, & \text{para } \dfrac{1}{4}\leq r<\dfrac{3}{4},\\[0.4cm]
            1, & \text{para } \dfrac{3}{4}\leq r.
        \end{cases}
        \]
        Determine la función de probabilidad de $X$.
    \end{mdframed}
\end{document}